\section{Функціональні вимоги}\label{sec:funktsional}

Усі функціональні вимоги системи мають стабільні унікальні ідентифікатори формату \texttt{REQ-F-xxx}, супроводжуються посиланнями на консенсусний простір знань $K^*$ та містять чіткі умови верифікації.

\subsection{Каталог фільмів і розклад сеансів}\label{sec:funktsional-katalog}

\begin{itemize}[leftmargin=*,labelindent=0pt]
    \item \textbf{REQ-F-001 (Загальнодоступний перегляд розкладу):} Система повинна надавати можливість перегляду актуального списку фільмів та сеансів будь-якому відвідувачеві без проходження автентифікації.~\sked{K2,K22}
    \textit{Умова верифікації:} Запит списку фільмів і сеансів без сесійного токена повертає успішний статус та коректні дані розкладу.
    
    \item \textbf{REQ-F-002 (Відображення схеми залу):} Система повинна відображати схему кінозалу обраного сеансу із актуальним статусом кожного місця («вільне» або «зайняте»). Вибір місця у вебінтерфейсі не резервує його і не блокує для інших користувачів.~\sked{K2,K31}
    \textit{Умова верифікації:} Вибір місця одним користувачем у браузері не змінює статус місця для інших користувачів до моменту підтвердження бронювання.
\end{itemize}

\subsection{Автентифікація, реєстрація та управління доступом}\label{sec:funktsional-auth}

\begin{itemize}[leftmargin=*,labelindent=0pt]
    \item \textbf{REQ-F-003 (Реєстрація та вхід за email):} Для створення та скасування бронювань користувач повинен зареєструватися й увійти в систему за допомогою дійсного email та пароля.~\sked{K23}
    \textit{Умова верифікації:} Спроба створення або скасування бронювання без автентифікації повертає статус відмови в авторизації.
    
    \item \textbf{REQ-F-004 (Внутрішній ID користувача):} Кожен зареєстрований обліковий запис отримує унікальний внутрішній ID, який пов'язується з усіма його транзакціями та логами.~\sked{K24}
    \textit{Умова верифікації:} Внутрішній ID користувача генерується автоматично, є незмінним та унікальним у сховищі даних.
    
    \item \textbf{REQ-F-005 (Ізоляція доступу до бронювань):} Користувач має право переглядати та скасовувати виключно власні бронювання.~\sked{K24}
    \textit{Умова верифікації:} Спроба перегляду або скасування бронювання користувача А за ідентифікатором користувача Б повертає статус заборони доступу (Forbidden).
    
    \item \textbf{REQ-F-006 (Відокремлення прав адміністратора):} Адміністратор працює в системі через окремий розділ керування довідниками. Звичайний користувач суворо позбавлений доступу до адміністративного розділу.~\sked{K26}
    \textit{Умова верифікації:} Запит користувача з роллю «Користувач» на адміністративний маршрут повертає статус заборони доступу (Forbidden).
\end{itemize}

\subsection{Створення бронювання та атомарне запобігання конфліктам}\label{sec:funktsional-bronuvannia}

\begin{figure}[htbp]
\centering
\begin{tikzpicture}[node distance=10mm and 12mm]
  \node[sked box] (free) {Вільне місце\\в залі};
  \node[sked box, right=of free] (booked) {Статус\\«Активне»};
  \node[sked box, right=of booked] (cancelled) {Статус\\«Скасоване»};
  \draw[sked arrow] (free) -- node[sked label, above=2pt] {Атомарне бронювання} (booked);
  \draw[sked arrow] (booked) -- node[sked label, above=2pt] {Скасування (до сеансу)} (cancelled);
  \draw[sked back] ([yshift=-4mm]cancelled.south) to[out=-135,in=-45] 
        node[sked label, below=2pt] {Місця знову вільні} ([yshift=-4mm]free.south);
\end{tikzpicture}
\caption{Життєвий цикл стану бронювання та місця}\label{fig:booking-lifecycle}
\end{figure}

\begin{itemize}[leftmargin=*,labelindent=0pt]
    \item \textbf{REQ-F-007 (Діапазон кількості місць):} Одне бронювання повинно містити від 1 до 6 місць включно.~\sked{K4}
    \textit{Умова верифікації:} Запит на створення бронювання з 0 або понад 6 місцями відхиляється з повідомленням про порушення обмеження діапазону.
    
    \item \textbf{REQ-F-008 (Атомарне створення без утримання):} Перевірка доступності всіх обраних місць та створення бронювання повинні виконуватися як єдина неподільна транзакція в момент натискання кнопки «Забронювати». Тимчасове утримання (hold) не застосовується.~\sked{K5}
    \textit{Умова верифікації:} У базі даних створення запису бронювання та маркування місць виконуються в одній транзакції без проміжних станів.
    
    \item \textbf{REQ-F-009 (Вирішення конкурентних конфліктів):} При одночасній спробі кількох користувачів забронювати одне й те саме місце на той самий сеанс успішним може бути лише одне бронювання. Усі інші запити відхиляються з повідомленням про зайнятість місця.~\sked{K29}
    \textit{Умова верифікації:} При імітації 10 паралельних запитів на одне місце успішним стає рівно один, а 9 запитів отримують відмову зі статусом конфлікту через зайнятість місця.
    
    \item \textbf{REQ-F-010 (Принцип «все або нічого»):} Якщо бронювання містить кілька місць і хоча б одне з них уже зайняте іншим користувачем, усе бронювання відхиляється повністю. Часткове бронювання суворо заборонене.~\sked{K30}
    \textit{Умова верифікації:} Запит на 3 місця, де 2 вільні й 1 зайняте, відхиляється без створення часткового бронювання; усі вільні місця залишаються доступними.
\end{itemize}

\subsection{Скасування бронювання користувачем}\label{sec:funktsional-skasuvannia-user}

\begin{itemize}[leftmargin=*,labelindent=0pt]
    \item \textbf{REQ-F-011 (Повне скасування бронювання):} Користувач може скасувати лише власне бронювання і виключно в повному обсязі (усі місця в ньому одночасно). Часткове скасування окремих місць заборонене.~\sked{K36}
    \textit{Умова верифікації:} Запит на скасування переводить усе бронювання у статус скасованого, вибіркові місця не скасовуються.
    
    \item \textbf{REQ-F-012 (Часовий дедлайн скасування):} Скасування бронювання дозволене виключно до настання часу початку сеансу ($T_{поточний} < T_{сеансу}$). У момент початку сеансу та після нього операція скасування блокується системою.~\sked{K37}
    \textit{Умова верифікації:} Запит на скасування, виконаний на 1 секунду пізніше за час початку сеансу, відхиляється з кодом помилки.
    
    \item \textbf{REQ-F-013 (Звільнення місць після скасування):} Після успішного скасування бронювання отримує статус «Скасоване», а всі зарезервовані за ним місця негайно стають доступними для вибору іншими користувачами.~\sked{K38}
    \textit{Умова верифікації:} Після успішного скасування повторний запит схеми залу відображає ці місця як вільні для нового бронювання.
    
    \item \textbf{REQ-F-014 (Ідемпотентність повторного скасування):} Повторний запит на скасування вже скасованого бронювання не змінює стан системи, не генерує побічних ефектів та є безпечним.~\sked{K39}
    \textit{Умова верифікації:} Повторний виклик скасування повертає інформацію про скасований статус без помилки та без повторного звільнення місць.
\end{itemize}

\subsection{Керування та скасування сеансів адміністратором}\label{sec:funktsional-admin}

\begin{itemize}[leftmargin=*,labelindent=0pt]
    \item \textbf{REQ-F-015 (Керування довідниками):} Адміністратор здійснює базові операції внесення та редагування інформації про фільми, кінозали та розклад сеансів.~\sked{K3}
    \textit{Умова верифікації:} Адміністратор може успішно створити новий сеанс для наявного фільму та кінозалу.
    
    \item \textbf{REQ-F-016 (Вільне редагування сеансів без бронювань):} Якщо на сеанс немає жодного активного бронювання, адміністратор має змогу вільно змінювати час початку, зал або фільм.~\sked{K44}
    \textit{Умова верифікації:} Редагування часу сеансу за відсутності бронювань завершується успішно з оновленням розкладу.
    
    \item \textbf{REQ-F-017 (Блокування зміни сеансів з бронюваннями):} Якщо на сеанс створено хоча б одне активне бронювання, пряме редагування параметрів сеансу (час, зал, фільм) повністю блокується системою.~\sked{K45}
    \textit{Умова верифікації:} Спроба оновлення часу або залу сеансу з активними бронюваннями повертає статус конфлікту через наявність активних бронювань.
    
    \item \textbf{REQ-F-018 (Явне скасування сеансу з причиною):} За наявності активних бронювань адміністратор має право явно скасувати сеанс за умови обов'язкового зазначення причини скасування та явного підтвердження дії.~\sked{K46}
    \textit{Умова верифікації:} Запит на скасування сеансу без заповненого поля причини відхиляється валідатором.
    
    \item \textbf{REQ-F-019 (Каскадне скасування бронювань):} При скасуванні сеансу сам сеанс та всі пов'язані з ним активні бронювання отримують статус «Скасовано» й зберігаються в історії для аудиту.~\sked{K47}
    \textit{Умова верифікації:} Після скасування сеансу всі бронювання цього сеансу в базі даних автоматично мають статус «Скасовано».
    
    \item \textbf{REQ-F-020 (Відображення причини в кабінеті користувача):} Користувач повинен бачити факт скасування свого бронювання адміністратором та текстову причину скасування сеансу у власному особистому кабінеті.~\sked{K48}
    \textit{Умова верифікації:} У відповіді API списку бронювань користувача для скасованого сеансу повертається поле з причиною від адміністратора.
\end{itemize}
